\chapter{Economy}
\section{Pillar 1: Money, Banking, Finance \& Insurance}
\section{Pillar 2: Budget, Taxation \& Public Finance}

\section{Pillar 3: Balance of Payments (BoP) \& International Trade}



\section{Pillar 4: Sectors of Economy, GDP \& Inflation}

\subsection{Union Budget 2026--27: Key Manufacturing Hub Announcements}

\begin{itemize}

  \item \textbf{Chemicals}: 3 new \textbf{Chemical Parks} -- cluster-based manufacturing; shared utilities, safety infrastructure, logistics connectivity.

  \item \textbf{Textiles}: Continued rollout of \textbf{7 PM MITRA Parks} (Mega Integrated Textile Region and Apparel) -- large-scale integrated textile manufacturing hubs.

  \item \textbf{MSME Manufacturing Clusters}: Strengthening cluster-based manufacturing + \textbf{Common Facility Centres (CFCs)} -- shared production, testing, quality infrastructure.

  \item \textbf{Capital Goods / Industrial Support}: Support for \textbf{tool rooms, testing \& calibration facilities} -- boost domestic capital goods manufacturing capability.

  \item \textbf{Infrastructure (PM GatiShakti)}: Continued multimodal, coordinated infrastructure planning linked to industrial/manufacturing locations.

  \item \textbf{Biopharma SHAKTI} (Global Biopharma \& Manufacturing Hub):
    \begin{itemize}
      \item Outlay: \textbf{₹10,000 crore / 5 years}.
      \item 3 new \textbf{NIPERs} (National Institutes of Pharmaceutical Education and Research).
      \item \textbf{1,000+ accredited clinical trial sites} (nationwide network).
      \item Objective: Position India as global biopharma manufacturing hub.
    \end{itemize}

\end{itemize}

\section{Pillar 5: Infrastructure}
\subsection{The Great Nicobar Project: Strategic, Economic and Ecological Framework}

\begin{itemize}
    \item \textbf{Context \& Strategic Significance:} 
    \begin{enumerate}
        \item Aimed at transforming Great Nicobar into a strategic maritime and economic hub in the Indian Ocean Region, strengthening India's presence in the Andaman Sea and Southeast Asia.
        \item Seeks to capture transshipment trade currently lost to foreign ports like Colombo, Singapore, and Klang due to India's lack of deep-water berths.
    \end{enumerate}

    \item \textbf{Core Infrastructure Components:}
    \begin{enumerate}
        \item \textbf{International Container Transhipment Terminal (ICTT):} Located at Galathea Bay with a capacity of 14.2 million Twenty-Foot Equivalent Units (TEU).
        \item \textbf{Greenfield International Airport:} Designed to handle 4,000 Peak Hour Passengers (PHP), projecting an initial 1 million passengers annually, scaling up to 10 million.
        \item \textbf{Township and Area Development:} A planned urban center spanning 16,610 hectares to support port-led residential, commercial, and institutional requirements.
        \item \textbf{Power Plant:} A 450 MVA hybrid gas and solar-based power plant designed with N-1 contingency (redundancy) to phase out diesel generators and reduce emissions.
    \end{enumerate}

    \item \textbf{Forest Diversion \& Compensatory Afforestation:}
    \begin{enumerate}
        \item Diverts 1.82\% of the total forest cover of the A\&N Islands.
        \item Preservation of 65.99 sq. km of land as strict green zones within the project area.
    \end{enumerate}

    \item \textbf{Tribal Welfare \& Rights Protection:}
    \begin{enumerate}
        \item Inhabited by Particular Vulnerable Tribal Groups (PVTGs): Shompens ($\approx$237, hunter-gatherers) and Nicobarese ($\approx$1,094, coastal fishing communities).
        \item \textbf{No Displacement/Relocation:} Project aligns with the Shompen Policy (2015), Jarawa Policy (2004), and Article 338A(9) of the Constitution. No tribal habitations are being relocated.
        \item \textbf{Tribal Reserve Land Balance:} 84.10 sq. km overlaps with the current Tribal Reserve. Within this, 11.032 sq. km is already utilized as revenue land since 1972. Thus, 73.07 sq. km is being de-notified for the project. To balance this, 76.98 sq. km is being re-notified as Tribal Reserve, creating a \textbf{net addition of 3.912 sq. km} to the reserve area.
    \end{enumerate}
\end{itemize}

\section{Pillar 6: Human Resource Development (HRD)}

\subsection{Women-Led Development: Welfare to Empowerment (2014--2026)}

\textbf{Paradigm Shift:} Welfare beneficiary $\rightarrow$ Women-led development; lifecycle approach across health, education, skills, finance, safety, governance.

\subsubsection*{A. Maternal \& Child Health}

\textbf{Beti Bachao Beti Padhao (BBBP)} [launched 22 Jan 2015]:
\begin{itemize}
  \item Objectives: Address declining CSR; enforce PCPNDT Act; promote girl child survival, protection, education; menstrual hygiene awareness.
  \item Impact: Sex ratio improved -- NFHS-5: \textbf{1,020 women/1,000 men} (vs.\ 943 in Census 2011).
\end{itemize}

\textbf{Pradhan Mantri Matru Vandana Yojana (PMMVY)} [2017]:
\begin{itemize}
  \item Maternity benefit: \textbf{\rupee5,000} (first child, 2 instalments); \textbf{\rupee6,000} (second child, girl only); DBT-linked to health milestones.
\end{itemize}

\textbf{Pradhan Mantri Surakshit Matritva Abhiyan (PMSMA)} [2016]:
\begin{itemize}
  \item Free ANC check-ups on \textbf{9th of every month} at government facilities; 2nd/3rd trimester focus.
  \item MMR declined: \textbf{88/lakh live births (2021--23)} from 130 (2014--15).
\end{itemize}

\textbf{Janani Suraksha Yojana (JSY)} [under NHM]:
\begin{itemize}
  \item Conditional cash for facility-based delivery; targets BPL, SC/ST women; ASHA workers as link.
\end{itemize}

\textbf{Janani Shishu Suraksha Karyakram (JSSK)} [extended 2014]:
\begin{itemize}
  \item Free entitlements for antenatal/post-natal complications; covers sick newborns up to 1 year.
\end{itemize}

\textbf{Key ANC/Delivery Outcomes (NFHS-4 $\rightarrow$ NFHS-6):}
\begin{itemize}
  \item Institutional births: 79\% $\rightarrow$ 90.6\%
\end{itemize}

\subsubsection*{B. Education}

\textbf{NEP 2020}: Gender Inclusion Fund; flexible learning pathways; multidisciplinary options.

\textbf{Samagra Shiksha} [2018--19]:
\begin{itemize}
  \item Integrated Pre-primary to Class XII; female enrolment: 1.57 crore (32\%, 2014--15) $\rightarrow$ \textbf{11.93 crore (48\%, 2024--25)}.
  \item Infrastructure (2024--25): 99.3\% schools with drinking water; 97.3\% with girls' toilets; 93.6\% with electricity.
\end{itemize}

\textbf{Kasturba Gandhi Balika Vidyalayas (KGBVs)} [under Samagra Shiksha]:
\begin{itemize}
  \item Residential schools (Class VI--XII) for SC/ST/OBC/disadvantaged girls in educationally backward blocks.
\end{itemize}

\textbf{Scholarships:}
\begin{itemize}
  \item Central Sector Scholarship (College/University): 50\% reserved for girls.
  \item AICTE PRAGATI Scholarship: 10,000/year since 2014--15; \textbf{36,000 girl students} by 2024--25.
\end{itemize}

\textbf{STEM:}
\begin{itemize}
  \item Vigyan Jyoti Scheme (Dec 2019): Classes IX--XII STEM mentoring; \textbf{1.12 lakh girls} across 300 districts, 34 States/UTs.
  \item IIT/NIT supernumerary seats: Women share $<$10\% $\rightarrow$ $>$20\%.
  \item UGC NET-JRF STEM (2024--25): Women -- \textbf{53\%+}.
\end{itemize}

\subsubsection*{C. Skills \& Digital Inclusion}

\textbf{PMKVY} [2015]:
\begin{itemize}
  \item PMKVY 4.0 focus: AI, drones, green energy, electronics; OJT component.
\end{itemize}

\textbf{NAVYA} (Nurturing Aspirations through Vocational Training for Young Adolescent Girls) [2025, under PMKVY 4.0]:
\begin{itemize}
  \item Target: Girls 16--18 yrs; digital marketing, cybersecurity, AI, green jobs, financial literacy.
\end{itemize}

\subsubsection*{D. Health \& Nutrition}

\textbf{Ayushman Bharat (2018) -- Women Dimensions:}
\begin{itemize}
  \item AB-PMJAY: 43.52 crore Ayushman cards (Feb 2026); women -- \textbf{49\% (21 crore)}; 4.97 crore women -- authorised hospital admissions; 36,229 empanelled hospitals.
  \item Ayushman Arogya Mandirs (AAM): \textbf{1.84 lakh} operational (Feb 2026); maternal, NCD, reproductive healthcare.
  \item ABDM: \textbf{95 crore health records} linked digitally; women -- \textbf{49.75\%} of ABHA accounts.
\end{itemize}

\textbf{Saksham Anganwadi \& Poshan 2.0:}
\begin{itemize}
  \item Targets: Children 0--6 yrs, adolescent girls, pregnant/lactating mothers..
\end{itemize}

\textbf{Mission Indradhanush} [Dec 2014]:
\begin{itemize}
  \item U-WIN digital tracking; 11.87 crore children + \textbf{3.96 crore pregnant women} registered; 1.32 crore pregnant women vaccinated; \textbf{8.73 crore} women screened for cervical cancer (as on Feb 2026).
\end{itemize}

\subsubsection*{E. Financial Inclusion \& Economic Empowerment}

\textbf{Sukanya Samriddhi Yojana (SSY)} [2015, under BBBP]:
\begin{itemize}
  \item Account for girls up to 10 yrs; deposits from \rupee250; interest \textbf{8.2\% p.a.}; tax-free (Sec 80C); partial withdrawal for education/marriage.
\end{itemize}

\textbf{DAY-NRLM / SHG Ecosystem:}
\begin{itemize}
  \item Coverage: 7,627 blocks; 1.51 crore community cadre; \textbf{93.85 lakh SHGs}; \textbf{10.07 crore members}.
  \item Bank credit accessed: \textbf{\rupee12.18 lakh crore}; 50,548 Bank Sakhis; 5.88 lakh enterprises under SVEP.
\end{itemize}

\textbf{Lakhpati Didi Scheme:}
\begin{itemize}
  \item Target: Sustainable annual income $\geq$ \rupee1 lakh/SHG woman.
  \item Government target: \textbf{6 crore Lakhpati Didis}.
\end{itemize}

\textbf{PMJDY} [2014]: Zero-balance accounts; foundation for DBT, insurance, credit access.

\textbf{Pradhan Mantri Mudra Yojana (PMMY)} [2015]:
\begin{itemize}
  \item Collateral-free loans: Shishu, Kishor, Tarun, Tarun Plus categories.
\end{itemize}

\textbf{PM SVANidhi} [Jun 2020]:
\begin{itemize}
  \item Collateral-free WC loans for urban street vendors: \rupee15k $\rightarrow$ \rupee25k $\rightarrow$ \rupee50k (on timely repayment); 7\% interest subsidy; UPI/RuPay linkage.
  \item 1.15 crore loans; 74.9 lakh vendors; women -- \textbf{46\% of beneficiaries}.
\end{itemize}

\textbf{Stand-Up India} [Apr 2016 -- Mar 2025]:
\begin{itemize}
  \item Loans \rupee10 lakh -- \rupee1 crore for new enterprises; 7-yr repayment with moratorium.
  \item \textbf{2.05 lakh women entrepreneurs} supported; women accounts: 55k (2018) $\rightarrow$ 1.90 lakh (2024).
\end{itemize}

\textbf{NaMo Drone Didi Yojana} [Nov 2023]:
\begin{itemize}
  \item Drones to SHG women for agri services (fertilizer/pesticide spraying); outlay \textbf{\rupee1,261 crore (2023--26)}.
  \item Target: 15,000 drones to SHGs; 1,094 drones distributed (2023--24); training at DGCA-authorised RPTOs.
\end{itemize}

\textbf{Womaniya (GeM)} [2019]:
\begin{itemize}
  \item Women SHG/enterprise products on Government e-Marketplace; direct procurement by Ministries/PSUs.
\end{itemize}

\textbf{SHE-Mart} [announced Union Budget 2026--27]:
\begin{itemize}
  \item Dedicated retail spaces (community-owned, SHG federation-managed) for women to sell directly to consumers.
  \item Target: \textbf{1 crore women beneficiaries}.
\end{itemize}

\subsubsection*{F. Safety, Sanitation \& Living Standards}

\textbf{Mission Shakti} [from Apr 2022]:
\begin{itemize}
  \item Two verticals: \textbf{Sambal} (safety -- OSCs, Helpline 181, Nari Adalat) \& \textbf{Samarthya} (empowerment -- Shakti Sadan, Sakhi Niwas, Palna Creches, SANKALP/HEW).
  \item 973 One Stop Centres; \textbf{14.49 lakh women assisted}; \textbf{3 crore+} supported via Helpline 181; 15,000+ Women Help Desks at police stations; integrated with ERSS-112.
\end{itemize}

\textbf{Swachh Bharat Mission} [Oct 2014]:
\begin{itemize}
  \item ODF status achieved Oct 2019; 6.3 lakh Community/Public Toilets (Urban); 2.7 lakh Community Sanitary Complexes (Rural); 12 crore+ individual toilets; 5 lakh+ ODF Plus (Model) villages.
\end{itemize}

\textbf{PMAY (Gramin \& Urban):}
\begin{itemize}
  \item PMAY-G (as Mar 2026): 4.15 crore allocated; 3.90 crore sanctioned; \textbf{2.99 crore completed}; \rupee4.03 lakh crore disbursed.
  \item PMAY-U 2.0: \textbf{96\% houses allotted to women}; 98.10 lakh homes completed.
\end{itemize}

\textbf{PM Ujjwala Yojana (PMUY)} [May 2016]:
\begin{itemize}
  \item Deposit-free LPG to BPL women; refill subsidy support.
  \item \textbf{10.55 crore connections} released (as on 14 May 2026).
\end{itemize}

\textbf{Jal Jeevan Mission (JJM)} [Aug 2019, extended to 2028]:
\begin{itemize}
  \item Functional Household Tap Connections (Har Ghar Jal); \textbf{15.84 crore rural families} connected; reduces women's water-fetching burden.
\end{itemize}

\subsubsection*{G. Governance \& Political Representation}

\begin{itemize}
  \item Women electorate (2024): \textbf{48.62\%}; 47 crore registered voters; voter turnout: \textbf{65.78\%} (marginally higher than men).
  \item Lok Sabha (2024): \textbf{75 women MPs}; Rajya Sabha: $\sim$17\%.
  \item PRIs: \textbf{14.5 lakh+} women elected representatives ($\sim$\textbf{46\% of total}).
  \item \textbf{Nari Shakti Vandan Adhiniyam, 2023}: \textbf{33\% reservation} for women in Lok Sabha and State Legislative Assemblies.
  \item NDA (National Defence Academy): First batch of 17 women cadets (2025); total \textbf{158 women cadets} by early 2026 (women inducted since 2022).
\end{itemize}

\subsection{Skilling India for a Future-Ready Workforce (Budget 2026--27 \& Policy Framework)}

\subsubsection*{Context \& Macro Indicators}
\begin{itemize}
  \item Economic Survey 2025--26: employment-focused skilling bridges skill gaps, improves productivity, promotes social mobility.
  \item PLFS (Feb 2026): LFPR (15+ yrs) = \textbf{55.9\%}; rising female participation; declining unemployment (rural + urban).
  \item Union Budget 2026--27: skill development as \textbf{cross-sectoral priority}.
\end{itemize}

\subsubsection*{Budget 2026--27: Sector-wise Skilling Initiatives}

\paragraph{Education} Allocation: \textbf{₹1.39 lakh crore} (↑8.27\% over BE 2025--26).
\begin{itemize}
  \item 5 \textbf{University Townships} near industrial/logistics corridors (universities, research, skill centres).
  \item 1 \textbf{girls' hostel per district} via VGF/capital support (STEM retention).
  \item Upgradation/establishment of \textbf{4 telescope facilities}: National Large Solar Telescope, National Large Optical Infrared Telescope, Himalayan Chandra Telescope, COSMOS-2 Planetarium.
\end{itemize}

\paragraph{Textiles \& Apparel} Sector: $\sim$2\% GDP, 11\% manufacturing GVA, 8.63\% exports; 45 million direct jobs; 6th largest global exporter.
\begin{itemize}
  \item \textbf{National Fibre Scheme} -- self-reliance in natural/man-made/new-age fibres.
  \item \textbf{Textile Expansion and Employment Scheme} -- modernise clusters.
  \item \textbf{National Handloom and Handicraft Programme} -- integrate schemes, support weavers/artisans.
  \item \textbf{Tex-Eco Initiative} -- sustainable/globally competitive textiles.
  \item \textbf{Samarth 2.0} -- industry-academia textile skilling.
  \item \textbf{Mega Textile Parks} -- scale + value addition in technical textiles.
  \item \textbf{Mahatma Gandhi Gram Swaraj Initiative} -- khadi/handloom/handicrafts branding + skilling.
\end{itemize}

\paragraph{Health} Outlay: \textbf{₹980 crore} (3 yrs) for allied/healthcare workforce.
\begin{itemize}
  \item \textbf{Biopharma SHAKTI}: ₹10,000 crore/5 yrs; 3 new NIPERs + upgradation of 7 existing NIPERs; 1,000+ clinical trial sites; strengthening CDSCO.
  \item Upgradation of \textbf{Allied Health Professionals (AHP) institutions}; 10 disciplines (optometry, radiology, anaesthesia, OT Tech, applied psychology etc.); \textbf{1 lakh AHPs added in 5 yrs}.
  \item \textbf{Care Ecosystem}: NSQF-aligned programmes; \textbf{1.5 lakh caregivers} trained in FY2026--27.
  \item \textbf{5 Regional Medical Hubs} (State-private partnership): AYUSH centres, Medical Value Tourism facilitation, diagnostics, rehab.
\end{itemize}

\paragraph{AYUSH} Allocation: \textbf{₹4,408 crore}.
\begin{itemize}
  \item 3 new \textbf{All India Institutes of Ayurveda (AIIA)}.
  \item Upgradation of AYUSH pharmacies and Drug Testing Laboratories.
  \item Upgradation of \textbf{WHO Global Traditional Medicine Centre}, Jamnagar.
\end{itemize}

\paragraph{Livestock} \textbf{₹6,153 crore}; contributes $\sim$16\% to farm income.
\begin{itemize}
  \item \textbf{Loan-linked capital subsidy}: private veterinary/para-vet colleges, hospitals, diagnostic labs; scale up veterinary professionals by \textbf{20,000+}.
\end{itemize}

\paragraph{Orange Economy (AVGC)} Media \& entertainment: ₹2.5 trillion (2024) → ₹3.06 trillion by 2027 (CAGR $\sim$7\%).
\begin{itemize}
  \item \textbf{AVGC Content Creator Labs}: 15,000 secondary schools + 500 colleges.
  \item AVGC sector demand: \textbf{2 million professionals by 2030}.
\end{itemize}

\paragraph{Design}
\begin{itemize}
  \item New \textbf{National Institute of Design} in eastern India.
\end{itemize}

\paragraph{Tourism} Sector: 5.22\% GDP (FY2024); \textbf{8.46 crore jobs} (13.3\% total employment).
\begin{itemize}
  \item \textbf{National Institute of Hospitality} -- upgrade of National Council for Hotel Management \& Catering Technology (NCHMCT); academia-industry-government bridge.
  \item Pilot: \textbf{upskilling 10,000 tourist guides} at 20 iconic sites (12-week hybrid course; IIM collaboration).
  \item \textbf{National Destination Digital Knowledge Grid} -- digital documentation of cultural/heritage sites.
\end{itemize}

\paragraph{Sports} Ministry allocation: ₹3,346 cr (RE 2025--26) → \textbf{₹4,479.88 cr} (BE 2026--27) (↑₹1,133 crore).
\begin{itemize}
  \item Target: \textbf{Top 10 by 2036; Top 5 by 2047}.
  \item \textbf{Khelo India Mission}: integrated talent pathway (foundational → elite), coach development, sports science integration, leagues/competitions, infrastructure.
\end{itemize}

\subsubsection*{Flagship Skilling Schemes (Skill India Mission, launched 2015)}

\paragraph{PMKVY / PMKVY 4.0}
\begin{itemize}
  \item Phases I--IV: pilot incentive → large-scale demand-driven skilling.
  \item \textbf{27.24 lakh trained} (38 sectors, 36 states, 737 districts) as of March 2026.
  \item April 2024--March 2026: \textbf{10.91 lakh trained} (IT-ITeS, aerospace, agri, rubber, leather etc.).
  \item \textbf{69 customised courses + 154 future-skill job roles} (AI, Industry 4.0, green jobs, digital).
  \item \textbf{16,900+ institutions}; 6,800+ Skill Hubs in schools/HEIs/ITIs.
\end{itemize}

\paragraph{Jan Shikshan Sansthan (JSS)} Origin: Shramik Vidyapeeth, 1967; 100\% GoI grant via NGOs.
\begin{itemize}
  \item Target: non-literates, neo-literates, school dropouts.
  \item \textbf{36,48,692 beneficiaries} trained (2018--March 2026); 26,519 tribal beneficiaries.
  \item 32 new NCVET-approved courses (NSQF Levels 3, 3.5, 4).
  \item Products marketed on \textbf{UdyamKart portal} (from Dec 2024).
\end{itemize}

\paragraph{NAPS / NAPS 2.0} Launched August 2016; ``earn while you learn'' model.
\begin{itemize}
  \item GoI contributes \textbf{25\% stipend ($\le$ ₹1,500/month)} via DBT.
  \item \textbf{54.41 lakh+ apprentices} engaged (2016--March 2026).
  \item FY2025--26: \textbf{12.35 lakh engaged}; 6.42 lakh completed OJT.
  \item \textbf{Certificate of Proficiency (CoP)}: launched Sep 2025; 67,711 CoPs generated by March 2026.
  \item FY2025--26 DBT: \textbf{₹562.75 crore} via 40.10 lakh transactions.
\end{itemize}

\paragraph{Craftsmen Training Scheme (CTS)} Introduced \textbf{1950}.
\begin{itemize}
  \item Training in \textbf{169 courses} via \textbf{14,688 ITIs} (3,345 Govt + 11,343 Pvt).
  \item 14 new CTS courses + 22 revised (3 years); enrolment ↑ 12.51 lakh (FY23) → \textbf{14.70 lakh (FY26)}.
\end{itemize}

\subsubsection*{PM--SETU (Pradhan Mantri Skilling \& Employability Transformation through Upgraded ITIs)}
\begin{itemize}
  \item Launched: \textbf{October 2025}; Centrally Sponsored Scheme; outlay: \textbf{₹60,000 crore}.
  \item \textbf{Hub-and-spoke model}: 1,000 Govt ITIs (200 hub + 800 spoke).
  \item \textbf{Special Purpose Vehicles (SPVs)} with Anchor Industry Partners (AIPs) for co-ownership/co-management.
  \item \textbf{5 National Skill Training Institutes} (Bhubaneswar, Chennai, Hyderabad, Kanpur, Ludhiana) → National Centres of Excellence for Skilling (global partnerships).
  \item States co-creating upgradation: Haryana, Karnataka, Tamil Nadu.
\end{itemize}

\subsubsection*{Key Targets \& Vision}
\begin{itemize}
  \item \textbf{Viksit Bharat @2047}: human capital as core growth driver.
  \item Demand-driven training, measurable outcomes, inclusive access.
  \item Formalisation + productivity gains via healthcare, AVGC, tourism, care economy skilling.
\end{itemize}


\clearpage